\def\version{00000000}
\def\opts#1#2#3#4{({\bf A})\,{#1}\ \ ({\bf B})\,{#2}\ \ ({\bf C})\,{#3}\ \ ({\bf D})\,{#4}}
\documentclass[CJK]{article}
\usepackage{geometry,amssymb, amsmath}
\usepackage{CJKutf8}
\usepackage{amsmath}
\usepackage{amssymb}
\usepackage{graphicx}
\usepackage{eufrak}
\usepackage{color}
\usepackage{slashed}
\usepackage{simplewick}
\usepackage{tikz}
\usepackage{tcolorbox}
\graphicspath{{../figures/}}
%%figures
\def\lfig#1#2{\includegraphics[width=#1 in]{#2}}
\def\addfig#1#2{\begin{center}\includegraphics[width=#1 in]{#2}\end{center}}
\def\wulian{\includegraphics[width=0.18in]{emoji_wulian.jpg}}
\def\bigwulian{\includegraphics[width=0.35in]{emoji_wulian.jpg}}
\def\bye{\includegraphics[width=0.18in]{emoji_bye.jpg}}
\def\bigbye{\includegraphics[width=0.35in]{emoji_bye.jpg}}
\def\huaixiao{\includegraphics[width=0.18in]{emoji_huaixiao.jpg}}
\def\bighuaixiao{\includegraphics[width=0.35in]{emoji_huaixiao.jpg}}
\def\jianxiao{\includegraphics[width=0.18in]{emoji_jianxiao.jpg}}
\def\bigjianxiao{\includegraphics[width=0.35in]{emoji_jianxiao.jpg}}
%% colors
\def\blacktext#1{{\color{black}#1}}
\def\bluetext#1{{\color{blue}#1}}
\def\redtext#1{{\color{red}#1}}
\def\darkbluetext#1{{\color[rgb]{0,0.2,0.6}#1}}
\def\skybluetext#1{{\color[rgb]{0.2,0.7,1.}#1}}
\def\cyantext#1{{\color[rgb]{0.,0.5,0.5}#1}}
\def\greentext#1{{\color[rgb]{0,0.7,0.1}#1}}
\def\darkgray{\color[rgb]{0.2,0.2,0.2}}
\def\lightgray{\color[rgb]{0.6,0.6,0.6}}
\def\gray{\color[rgb]{0.4,0.4,0.4}}
\def\blue{\color{blue}}
\def\red{\color{red}}
\def\green{\color{green}}
\def\darkgreen{\color[rgb]{0,0.4,0.1}}
\def\darkblue{\color[rgb]{0,0.2,0.6}}
\def\skyblue{\color[rgb]{0.2,0.7,1.}}
%%control
\def\be{\begin{equation}}
\def\ee{\nonumber\end{equation}}
\def\bea{\begin{eqnarray}}
\def\eea{\nonumber\end{eqnarray}}
\def\bch{\begin{CJK}{UTF8}{gbsn}}
\def\ech{\end{CJK}}
\def\bitem{\begin{itemize}}
\def\eitem{\end{itemize}}
\def\bcenter{\begin{center}}
\def\ecenter{\end{center}}
\def\bex{\begin{minipage}{0.2\textwidth}\includegraphics[width=0.6in]{jugelizi.png}\end{minipage}\begin{minipage}{0.76\textwidth}}
\def\eex{\end{minipage}}
\def\chtitle#1{\frametitle{\bch#1\ech}}
\def\bmat#1{\left(\begin{array}{#1}}
\def\emat{\end{array}\right)}
\def\bcase#1{\left\{\begin{array}{#1}}
\def\ecase{\end{array}\right.}
\def\bmini#1{\begin{minipage}{#1\textwidth}}
\def\emini{\end{minipage}}
\def\tbox#1{\begin{tcolorbox}#1\end{tcolorbox}}
\def\pfrac#1#2#3{\left(\frac{\partial #1}{\partial #2}\right)_{#3}}
%%symbols
\def\bropt{\,(\ \ \ )}
\def\optlist#1#2#3{(A)\;#1\ \ \ (B)\;#2\ \ \ (C)\;#3}
\def\sone{$\star$}
\def\stwo{$\star\star$}
\def\sthree{$\star\star\star$}
\def\sfour{$\star\star\star\star$}
\def\sfive{$\star\star\star\star\star$}
\def\rint{{\int_\leftrightarrow}}
\def\roint{{\oint_\leftrightarrow}}
\def\stdHf{{\textit{\r H}_f}}
\def\deltaH{{\Delta \textit{\r H}}}
\def\ii{{\dot{\imath}}}
\def\skipline{{\vskip0.1in}}
\def\skiplines{{\vskip0.2in}}
\def\lagr{{\mathcal{L}}}
\def\hamil{{\mathcal{H}}}
\def\vecv{{\mathbf{v}}}
\def\vecx{{\mathbf{x}}}
\def\vecy{{\mathbf{y}}}
\def\veck{{\mathbf{k}}}
\def\vecp{{\mathbf{p}}}
\def\vecn{{\mathbf{n}}}
\def\vecA{{\mathbf{A}}}
\def\vecP{{\mathbf{P}}}
\def\vecsigma{{\mathbf{\sigma}}}
\def\hatJn{{\hat{J_\vecn}}}
\def\hatJx{{\hat{J_x}}}
\def\hatJy{{\hat{J_y}}}
\def\hatJz{{\hat{J_z}}}
\def\hatj#1{\hat{J_{#1}}}
\def\hatphi{{\hat{\phi}}}
\def\hatq{{\hat{q}}}
\def\hatpi{{\hat{\pi}}}
\def\vel{\upsilon}
\def\Dint{{\mathcal{D}}}
\def\adag{{\hat{a}^\dagger}}
\def\bdag{{\hat{b}^\dagger}}
\def\cdag{{\hat{c}^\dagger}}
\def\ddag{{\hat{d}^\dagger}}
\def\hata{{\hat{a}}}
\def\hatb{{\hat{b}}}
\def\hatc{{\hat{c}}}
\def\hatd{{\hat{d}}}
\def\hatN{{\hat{N}}}
\def\hatH{{\hat{H}}}
\def\hatp{{\hat{p}}}
\def\Fup{{F^{\mu\nu}}}
\def\Fdown{{F_{\mu\nu}}}
\def\newl{\nonumber \\}
\def\vece{\mathrm{e}}
\def\calM{{\mathcal{M}}}
\def\calT{{\mathcal{T}}}
\def\calR{{\mathcal{R}}}
\def\barpsi{\bar{\psi}}
\def\baru{\bar{u}}
\def\barv{\bar{\upsilon}}
\def\qeq{\stackrel{?}{=}}
\def\torder#1{\mathcal{T}\left(#1\right)}
\def\rorder#1{\mathcal{R}\left(#1\right)}
\def\contr#1#2{\contraction{}{#1}{}{#2}#1#2}
\def\trof#1{\mathrm{Tr}\left(#1\right)}
\def\trace{\mathrm{Tr}}
\def\comm#1{\ \ \ \left(\mathrm{used}\ #1\right)}
\def\tcomm#1{\ \ \ (\text{#1})}
\def\slp{\slashed{p}}
\def\slk{\slashed{k}}
\def\calp{{\mathfrak{p}}}
\def\veccalp{\mathbf{\mathfrak{p}}}
\def\Tthree{T_{\tiny \textcircled{3}}}
\def\pthree{p_{\tiny \textcircled{3}}}
\def\dbar{{\,\mathchar'26\mkern-12mu d}}
\def\erf{\mathrm{erf}}
\def\const{\mathrm{constant}}
\def\pheat{\pfrac p{\ln T}V}
\def\vheat{\pfrac V{\ln T}p}
%%units
\def\fdeg{{^\circ \mathrm{F}}}
\def\cdeg{^\circ \mathrm{C}}
\def\atm{\,\mathrm{atm}}
\def\angstrom{\,\text{\AA}}
\def\SIL{\,\mathrm{L}}
\def\SIkm{\,\mathrm{km}}
\def\SIyr{\,\mathrm{yr}}
\def\SIGyr{\,\mathrm{Gyr}}
\def\SIV{\,\mathrm{V}}
\def\SImV{\,\mathrm{mV}}
\def\SIeV{\,\mathrm{eV}}
\def\SIkeV{\,\mathrm{keV}}
\def\SIMeV{\,\mathrm{MeV}}
\def\SIGeV{\,\mathrm{GeV}}
\def\SIcal{\,\mathrm{cal}}
\def\SIkcal{\,\mathrm{kcal}}
\def\SImol{\,\mathrm{mol}}
\def\SIN{\,\mathrm{N}}
\def\SIHz{\,\mathrm{Hz}}
\def\SIm{\,\mathrm{m}}
\def\SIcm{\,\mathrm{cm}}
\def\SIfm{\,\mathrm{fm}}
\def\SImm{\,\mathrm{mm}}
\def\SInm{\,\mathrm{nm}}
\def\SImum{\,\mathrm{\mu m}}
\def\SIJ{\,\mathrm{J}}
\def\SIW{\,\mathrm{W}}
\def\SIkJ{\,\mathrm{kJ}}
\def\SIs{\,\mathrm{s}}
\def\SIkg{\,\mathrm{kg}}
\def\SIg{\,\mathrm{g}}
\def\SIK{\,\mathrm{K}}
\def\SImmHg{\,\mathrm{mmHg}}
\def\SIPa{\,\mathrm{Pa}}

\geometry{tmargin=0.3in, bmargin=0.4in, lmargin=0.7in, rmargin=0.8in, nohead, nofoot}
\begin{document}
\bch
\bcenter
第二届华南地区大学生天文竞赛 总赛 （共4页，满分100分) 

{\vskip 0.05in}

姓名 \uline{0.8} {\hskip 0.2in}  学校 \uline{1.5}{\hskip 0.2in} 学号 \uline{0.8}{\hskip 0.2in} 得分 \uline{0.5}

\ecenter

{\bf \noindent（一）选择题，每小题2分，共40分。请集中把答案写在下面的答案区:}

{\vskip 0.04in}

{(1) \uline{0.3}\ (2) \uline{0.3}\ (3) \uline{0.3}\ (4) \uline{0.3}\ (5) \uline{0.3}\ (6) \uline{0.3}\ (7) \uline{0.3}\ (8) \uline{0.3}\ (9) \uline{0.3}\ (10) \uline{0.23}}

{\vskip 0.04in}

{(11) \uline{0.23}\  (12) \uline{0.23}\ (13) \uline{0.23}\ (14) \uline{0.23}\ (15) \uline{0.23}\ (16) \uline{0.23}\ (17) \uline{0.23}\ (18) \uline{0.23}\ (19) \uline{0.23}\ (20) \uline{0.23}}

\bitem
\item[(1)]{日、月食只可能发生在哪些月相发生的时候

    \opts{望和朔}{上弦月和下弦月}{上弦月和望}{下弦月和望}
}    
\item[(2)]{下列哪个不属于类地行星？
  
  \opts{地球}{土星}{水星}{火星}
}
\item[(3)]{北极星属于哪个星座？

  \opts{大熊座}{小熊座}{双子座}{仙后座}
}
\item[(4)]{太阳黑子活动的周期大约是
  
  \opts{7年}{11年}{19年}{25年}
}
\item[(5)]{《春秋》记载“秋七月，有星孛入北斗”是世界上最早的关于\uline{0.5}的最早记录？

    \opts{哈雷彗星}{超新星爆发}{$\gamma$射线暴}{类星体}
}  
\item[(6)]{《周髀算经》是中国最古老的天文学和数学著作，一般认为它约成书于

  \opts{夏朝}{西周时期}{战国前后}{初唐}
}
\item[(7)]{我国古历中干支纪年的周期是

  \opts{十年}{十二年}{六十年}{一百二十年}
}
\item[(8)]{下列哪一颗恒星目前{\bf 不在}主序阶段？

    \opts{太阳}{牛郎星}{织女星}{北极星}
}    
\item[(9)]{根据现有恒星演化理论，太阳大概在多久以后变成红巨星

    \opts{1亿年}{20亿年}{100亿年}{500亿年}
}  
\item[(10)]{天狼星A的伴星，也就是天狼星B，亮度比天狼星A差10个星等，表面温度比太阳还高。它是一颗

    \opts{白矮星}{中子星}{黑洞}{红巨星}
}
\item[(11)]{著名的蟹状星云(NGC1952)是

    \opts{超新星爆发遗迹}{遥远的河外星系}{黑洞吸积盘}{恒星周围电离的气体和尘埃}
}  
\item[(12)]{一般认为，脉冲星是

    \opts{高速旋转的白矮星}{高速旋转的中子星}{高速旋转的黑洞}{表面不断膨胀收缩的恒星}
}
\item[(13)]{大麦哲伦云和小麦哲伦云绕我们银河系旋转。按照哈勃星系分类法，它们属于

    \opts{椭圆星系}{漩涡星系}{棒旋星系}{不规则星系}
}  
\item[(14)]{在宇宙学标准模型里，宇宙晚期的加速膨胀由什么造成？

    \opts{暗物质}{暗能量}{引力透镜效应}{宇宙微波背景辐射}
}  
\item[(15)]{可观测宇宙中，氦元素占所有元素的质量百分比大约为

    \opts{5\%}{25\%}{75\%}{95\%}
}
\item[(16)]{宇宙学标准模型预言了，但是目前还没有测量到的是

    \opts{宇宙微波背景辐射}{宇宙中微子背景辐射}{弱引力透镜}{重子声波振荡}
}  
\item[(17)]{下列哪一个望远镜是光学望远镜？

    \opts{郭守敬望远镜(LAMOST)}{平方千米阵(SKA)}{慧眼(HXMT)}{天眼(FAST)}
}  
\item[(18)]{一个直径6.5米的光学望远镜的衍射极限分辨率大约为

    \opts{0.02角秒}{0.2角秒}{2角秒}{20角秒}
}      
\item[(19)]{天文观测中的“蒙气差”，是指

  \opts{大气折射使天体看起来在天空中位置更高}{大气折射使天体看起来在天空中位置更低 \\}{大气散射使天体视位置发生随机移动}{大气散射使天体看起来模糊}
}
\item[(20)]{下列哪个天文学的成就获得了诺贝尔奖？
  
  \opts{爱因斯坦的广义相对论}{哈勃发现宇宙膨胀 \\}{霍金的黑洞辐射理论}{福勒的恒星化学元素形成理论}
}
  
\eitem

{\noindent {\bf (二)}
《周髀算经》中记载了用髀（直立竿）晷（在正午时的影子长度）标定地理位置的办法：“周髀长{\bf 八尺}……故冬至日晷{\bf 丈三尺五寸}，夏至日晷{\bf 尺六寸}。”
 请由此估算当时的黄赤交角和测量地的纬度。{\blue (10分)}

{\vskip 4in}

\noindent {\bf(三)}
某行星的地质构造和平均密度都和地球相近。表面具有稳定的、温度和地球的大气相近的大气层，其主要成分是氦气。请根据这些信息（必要的话加上一些合理的假设）粗略地估算该行星的半径。 {\blue (10分)}

\newpage
}
姓名 \uline{0.8} {\hskip 0.2in}  学校 \uline{1.5}{\hskip 0.2in} 学号 \uline{0.8}{\hskip 0.2in}

{\noindent {\bf (四)}
把地球理想化为一个密度均匀的，半长轴(赤道半径)为 $6380\SIkm$，半短轴(南北极距离的一半)为 $6360\SIkm$ 的扁平状旋转椭球体 (由椭圆绕短轴旋转而成)。忽略潮汐和其他天体的影响，估算在赤道上的重力常数和在北极（或南极）的重力常数的相对差异。{\blue (20分)}
\newpage
}
{\noindent {\bf (五)}
量子力学的海森堡测不准原理告诉我们：无法同时完全确定粒子的位置$x$和动量$p$，位置和动量的不确定性的乘积有个下限：$\Delta x \Delta p \gtrsim h$，其中$h=6.63\times 10^{-34}\mathrm{J\cdot s}$是普朗克常量。

现在考虑一个质量为$M$的史瓦西黑洞。它的视界半径$R_H$是多大？{\blue (5分)}把该黑洞当成一个半径为$R_H$的球形封闭空间，假想该封闭空间区域内有一些测试光子，用海森堡测不准原理估算每个光子的最小可能动量的数量级。{\blue (5分)} 如果认为这些测试光子处于某种温度为$T$的热平衡状态，估算$T$的下限。 {\blue (5分)} 把黑洞的辐射当成黑体辐射，估算现在能观测到的来自早期宇宙的黑洞的质量下限。{\blue (5分)}

{\vskip 7.8in}
}


常数表：

万有引力常数 $G = 6.7\times 10^{-11}\mathrm{N\cdot m^2/kg^2}$

太阳质量 $M_\odot = 2.0\times 10^{30}\SIkg$

天文单位（日地平均距离）$1\mathrm{AU} = 1.5\times 10^{11}\SIm$

玻尔兹曼常数 $k_B=1.38\times 10^{-23}\mathrm{J/K}$
\bcenter
    {\scriptsize 试卷印刷号: \version}
\ecenter
\ech
\end{document}
